\section{Reproducibility note: saved generation settings can change decoding}
\label{app:genconfig}

Our inference code passes an explicit greedy configuration to \texttt{generate} (\texttt{do\_sample=False}, one beam).
In the library version we use (\texttt{transformers} 4.57.1), the model's saved \texttt{generation\_config.json} still overrides every field that the caller left at its global default, provided the file records a library version of 4.50 or later.
\texttt{do\_sample=False} is itself the global default, so it can be overridden.

The released \texttt{Qwen2.5-1.5B-Instruct} configuration specifies sampling (temperature 0.7, top-$p$ 0.8, top-$k$ 20, repetition penalty 1.1), but it records library version 4.37.0, so the override does not apply and decoding stays greedy.
Saving a fine-tuned model with \texttt{save\_pretrained} rewrites the recorded version to 4.57.1.
With the same inference code, the fine-tuned model is then decoded by sampling with these settings.
We verified both cases by computing the effective generation configuration in a weight-free (meta-device) instantiation of the model.

Comparing a greedy baseline against sampled fine-tuned models would confound the comparison.
We therefore copy the original \texttt{generation\_config.json} byte for byte into every evaluated checkpoint.
We also check the effective configuration before each evaluation.

\section{Training precision pilot}
\label{app:precision}

The first M1 run was trained on an NVIDIA T4.
The software check for BF16 support reported that the GPU supports BF16, although the T4 has no native BF16 arithmetic, so the run trained in BF16.
We kept this run as a pilot.
We then fixed the training precision of every controlled run to FP16, the precision used at inference, and moved all controlled runs to one GPU type.
BF16 is now selected only on GPUs with native support, and controlled runs request FP16 explicitly.

Table~\ref{tab:precision-setup} compares the pilot with the controlled FP16 run of M1.
Both use the same training data version, hyperparameters, seed and number of optimizer steps.
They differ in training precision, training GPU and generation GPU, so the comparison cannot separate the effect of precision from that of the hardware.
Table~\ref{tab:precision-metrics} lists the scores on the development split for both runs and the base model.
Scores come from one judge pass without a significance test.
The two runs produced identical answers on 63 of the 156 development items.
Against the pilot, the FP16 run's reference verdict was higher on 12 items, unchanged on 137 and lower on 7.
The final validation loss was 1.6828 for the pilot and 1.6721 for the FP16 run.
These numbers describe the two runs only; we do not draw conclusions about precision from them.

% generated by ktide.precision_report; do not edit by hand
\begin{table}[t]
\centering
\small
\begin{tabular}{lrr}
\toprule
 & M1 pilot (BF16) & M1 (FP16) \\
\midrule
Training GPU & Tesla T4 & RTX 5000 Ada Generation \\
Training precision & BF16 & FP16 \\
Training time (s) & 423.8 & 18.6 \\
Validation loss (epoch 3) & 1.6828 & 1.6721 \\
Non-finite steps / scale drops & not logged & 0 / 0 \\
Generation GPU & Tesla T4 & RTX 5000 Ada Generation \\
\bottomrule
\end{tabular}
\caption{Two M1 training runs with the same data, hyperparameters and seed. They differ in training precision and GPU. Identical answers on the 156 shared items: 63.}
\label{tab:precision-setup}
\end{table}

\begin{table}[t]
\centering
\small
\begin{tabular}{lrrr}
\toprule
Metric & M0 & M1 pilot (BF16) & M1 (FP16) \\
\midrule
Strict accuracy (correct only) & 0.051 & 0.038 & 0.045 \\
Reference score (partial = 0.5) & 0.218 & 0.218 & 0.234 \\
Judge total (1-5) & 3.18 & 3.35 & 3.48 \\
Judge accuracy (1-5) & 2.28 & 2.23 & 2.33 \\
Judge relevance (1-5) & 3.27 & 3.18 & 3.34 \\
Judge constraints (1-5) & 3.98 & 4.65 & 4.76 \\
30-word limit & 0.788 & 0.974 & 0.994 \\
1-2 sentences & 0.782 & 1.000 & 1.000 \\
Both length limits & 0.731 & 0.974 & 0.994 \\
English term preservation & 0.328 & 0.100 & 0.089 \\
Polite register & 0.962 & 0.000 & 0.000 \\
Source mention & 0.000 & 0.013 & 0.013 \\
Chinese characters & 0.013 & 0.000 & 0.000 \\
\bottomrule
\end{tabular}
\caption{Development split, one judge pass, no significance test. ``pending'' marks scores not yet computed.}
\label{tab:precision-metrics}
\end{table}

